\documentclass[10pt,a4paper]{amsart}
\usepackage[margin=19mm]{geometry}
\usepackage{amsmath,amssymb,mathtools}
\usepackage[scr=rsfs,cal=euler]{mathalfa}
\usepackage{xfrac}
\usepackage[unicode,hidelinks,bookmarks=false]{hyperref}
\newcommand{\ie}{i.e.\ }
\newcommand{\cf}{cf.\ }
\newcommand{\resp}{respectively\ }
\newcommand{\Subsection}{\subsection}
\providecommand{\nobreakdash}{\nobreak\hbox{-}}
\setlength{\emergencystretch}{2em}
\multlinegap0pt
\usepackage{bm}
\usepackage{bbm}
\usepackage{dsfont}
\usepackage{xspace}
\usepackage{cancel}
\usepackage{mathtools}
\usepackage{amsthm}
\makeatletter
\@ifundefined{teo}{\newtheorem{teo}{Theorem}}{}
\makeatother
\newcommand{\N}{\mathbb{N}}
\newcommand{\R}{\mathbb{R}}
\title[A PINNs approach for the computation of eigenvalues in ellip]{A PINNs approach for the computation of eigenvalues in elliptic problems}
\author{Julian Fernandez Bonder, Ariel M. Salort}
\date{Source: arXiv:2507.03126v2 (CC BY 4.0)}
\begin{document}
\maketitle

\noindent\emph{Source and licence.} A PINNs approach for the computation of eigenvalues in elliptic problems. Version arXiv:2507.03126v2 (CC BY 4.0), distributed under the Creative Commons Attribution 4.0 International licence, as recorded in the arXiv licence metadata for this exact version and shown on the abstract page https://arxiv.org/abs/2507.03126v2. Full rights notice: licenses/arxiv-2507.03126-CC-BY-4.0.txt.\par\smallskip

\noindent\textit{Editorial scope.} Counted unit: the Teo whose source label is \texttt{teo.min} in the article (source0.tex lines 815--821, source label \texttt{teo.min}), delivered together with the article's own complete original proof of it (source0.tex lines 823--864, one \texttt{proof} environment of 42 lines). No mathematics is written, repaired or re-derived: statement and proof are the article's own text line for line. Required context retained uncounted: Lambda (lines 155--157); conf.V (lines 247--252). Every definition, display and label of those lines is retained unchanged. Omitted from this package and not redistributed: 1--154, 158--246, 253--814, 822--822, 865--865. Nothing inside a retained excerpt is omitted or altered: every retained body is delivered exactly as the article's lines are written, so no omission receipt of any kind accompanies this package. Citations: the retained excerpts carry no citation command at all, so no bibliography entry of the article is transcribed and no reference is redistributed. Reference adaptation: none was needed and none was made; every reference inside the delivered unit resolves inside it, so no unresolved reference or citation is produced. Printed slips kept literal: no printed word, symbol, hypothesis or number of the counted statement or of its proof was altered. Layout and class: the article's own class is \texttt{revtex4-2} with packages including xcolor, inputenc, bm, graphicx, booktabs, amsmath, amsthm, amssymb, amsfonts, mathtools, physics, enumitem, float, tikz; the wrapper is the lane's pinned \texttt{amsart} editorial wrapper at 10pt on A4, which restates the theorem-like environments the retained text needs and adds the article's own macro definitions that the retained text needs (\texttt{\textbackslash N}, \texttt{\textbackslash R}), with extra wrapper packages (bm, bbm, dsfont, xspace, cancel, mathtools). The delivered numbering is the unit's own. Assets: no retained span contains a figure environment, a graphics-inclusion command, a PDF-page inclusion, a table or a TikZ picture; no image file is copied into this package, the package declares no asset dependency and no image file is redistributed with it. Licence: version arXiv:2507.03126v2 of this article is distributed under the Creative Commons Attribution 4.0 International licence, as recorded in the arXiv licence metadata for this exact version and shown on the abstract page https://arxiv.org/abs/2507.03126v2. A full rights notice is delivered at licenses/arxiv-2507.03126-CC-BY-4.0.txt. Characters: every retained excerpt is pure ASCII, so the delivered engine, XeLaTeX with its default Latin Modern text font, reports no missing character.
\begin{equation}\label{Lambda}
\Lambda(u, E) = \frac{\int (\Delta u - V(x) u + Eu)^2\, dx}{\int u^2\, dx}.
\end{equation}

\begin{equation}\label{conf.V}
V(x) = \begin{cases}
0 & \text{if } x\in \Omega\\
\infty & \text{if } x\not\in \Omega,
\end{cases}
\end{equation}

\begin{teo}\label{teo.min}
Let $V(x)$ be as in \eqref{conf.V} with $\Omega\subset \R^n$ bounded. Given $E\ge 0$, there exists $u_E\in H^2(\Omega)\cap H^1_0(\Omega)$ such that
$$
\Lambda(E) = \Lambda(u_E, E) = \inf_{u\in H^2(\Omega)\cap H^1_0(\Omega)}\Lambda(u, E),
$$
where $\Lambda(u, E)$ is given in \eqref{Lambda}.
\end{teo}

\begin{proof}
The result follows from the \emph{direct method in the Calculus of Variations}. Fix $E \ge 0$, and consider a minimizing sequence $\{u_j\}_{j \in \N} \subset H^2(\Omega) \cap H^1_0(\Omega)$ such that $\|u_j\|_2 = 1$ and
$$
\Lambda(E) = \lim_{j \to \infty} \Lambda(u_j, E).
$$
In particular, the sequence $\Lambda(u_j, E)$ is bounded and nonnegative:
\begin{equation}\label{cota.Lambda}
0 \le \Lambda(u_j, E) \le C,
\end{equation}
for some constant $C > 0$.

Now observe that
\begin{align*}
\Lambda(u_j, E) &= \int_\Omega (\Delta u_j + Eu_j)^2\, dx \\
&= \int_\Omega (\Delta u_j)^2\, dx + 2E\int_\Omega u_j \Delta u_j \, dx+ E^2.
\end{align*}
Using the inequality $ab\le \epsilon a^2 + \frac{b^2}{4\epsilon}$ for any $a, b>0$ and $\epsilon >0$we estimate:
$$
\int_\Omega |u_j| |\Delta u_j|\, dx \le \epsilon \int_\Omega (\Delta u_j)^2\, dx + \frac{1}{4\epsilon}\int_\Omega u_j^2\, dx.
$$
Taking $\epsilon=1/4E$, we obtain
\begin{equation}\label{cota.Lambda2}
\Lambda(u_j, E)\ge \frac12 \int_\Omega (\Delta u_j)^2\, dx + E^2-E.
\end{equation}
Combining \eqref{cota.Lambda}, \eqref{cota.Lambda2} and the normalization $\|u_j\|_2=1$ we conclude that the sequence $\{u_j\}_{j\in\N}$ is bounded in $H^2(\Omega)\cap H^1_0(\Omega)$.

Therefore, up to a subsequence, we have $u_j \rightharpoonup u_E$ weakly in $H^2(\Omega) \cap H^1_0(\Omega)$ for some $u_E$ in that space.

Since the functional 
$$
u \mapsto \int_\Omega (\Delta u + E u)^2\, dx
$$ 
is convex in $u$, it is weakly lower semicontinuous. Thus,
$$
\int_\Omega (\Delta u_E + E u_E)^2\, dx \le \liminf_{j \to \infty} \int_\Omega (\Delta u_j + E u_j)^2\, dx.
$$
Moreover, the immersion $H^2(\Omega)\cap H^1_0(\Omega)\hookrightarrow L^2(\Omega)$ is compact, so
$$
\int_\Omega u_E^2\, dx = \lim_{j\to\infty} \int_\Omega u_j^2\, dx = 1.
$$
Hence, $u_E$ is admissible and attains the minimum.
\end{proof}

\end{document}
